\section{Least squares on magnitude images}\label{sec:magnitude}

The published inverse works with magnitudes, made real-valued by the sign convention of its
Eq.~18. A direct alternative is to fit the model magnitudes,
\begin{equation}
  \hat\theta=\arg\min_\theta\sum_j\bigl(|y_j|-|s_j(\theta)|\bigr)^2,
  \label{eq:magls}
\end{equation}
over $\theta=(\log\Mzero,\log\Bone,\log\Tone,\log\Ttwo,\log\Ttwos)$ with the same projected LM iteration,
constraints and four starts as complex ML. This is not maximum likelihood for magnitude data,
which are Rician: it ignores the noise floor $\mathbb E|y|-A\approx\sigma^2/(2A)$, which for the weakest
amplitude of the reference tissue is $10\%$ at SNR 300, $1\%$ at SNR 1000 and $0.1\%$ at SNR 3000.
\Cref{sec:rician} showed that magnitudes carry essentially all the local information; this
section concerns the global bit they lose.

\subsection{Why a magnitude fit needs the branch test and complex ML does not}

\paragraph{1. Magnitudes lose the sign of $\sin\alpha$.} By \cref{thm:uv} and \cref{lem:absv}
the magnitudes of scan $i$ determine $u_i$ and $\Mzero|v_i|=\Mzero\sqrt{1-u_i^2}\tanh(\TR/2\Tone)$, i.e.\
$\cos\alpha_i$ for any trial $E_1$, but not the sign of $\sin\alpha_i$. Equivalently, $|a_k(\alpha)|$ is
even and $2\pi$-periodic in $\alpha$, hence symmetric about every multiple of $180^\circ$.

\paragraph{2. This produces an exact alias in the physiological range.} By
\cref{prop:equiv}(a,c) the magnitude data at $(\Mzero,\Bone,\Tone)$ are reproduced exactly at every
solution $(B_1^{+\prime},E_1')$ of \eqref{eq:aliaseq}, with $\Mzero'$ set by part (a). The branch
$\alpha_2=360^\circ+x'$ admits such a solution: for the reference tissue $(1,1500\,\mathrm{ms},1)$ and
$(1.199,1038\,\mathrm{ms},0.832)$ give the same 18 magnitudes to machine precision
(\cref{tab:aliases}). The alias is close to the small-angle scaling
($1500/1.199^2=1043$) but exact.

\paragraph{3. It is a symmetry of the magnitude likelihood, not a noise effect.} The alias
exists for every parameter value, so the map $\theta\mapsto|s(\theta)|$ takes the same values on two
separate sheets of parameter space. Any likelihood that depends on the data only through
magnitudes, Gaussian or Rician, has two equal global maxima for every noise realisation. We
verified this directly: fitting noisy magnitudes on each side of $360^\circ$ gives final costs
that agree to a median relative difference of $10^{-12}$ at SNR 300, 1000 and 3000. No SNR
resolves it, and the local CRLB, identical on both sheets, cannot detect it. An
unconstrained magnitude fit returns the maximum nearest its starting point: in
\cref{tab:mag} it chose the wrong branch in 47, 3 and \MagWrongThree{} of 1000 voxels at SNR
300, 1000 and 3000, and in 13, 3 and 0 when the same fits were run with the exact Jacobian,
which takes slightly different paths. The counts describe the optimiser, not the data.

\paragraph{4. Complex data carry the missing bit.} By \cref{prop:period},
$a_k(-\alpha)=-a_k(\alpha)$, so moving $\alpha_2$ from $360^\circ-x$ to $360^\circ+x'$ flips the sign of every
scan-2 signal relative to scan 1 (\cref{prop:equiv}(b)). A single $\phi_0$ and a single $\dw$ are
shared by both scans, so the relative sign cannot be absorbed. The best complex fit
restricted to $360^\circ<\alpha_2<540^\circ$ moves to the edge $\alpha_2=360^\circ$, where scan 2 vanishes, and
leaves the entire scan-2 signal ($0.144\,\Mzero$) as residual: its cost exceeds the true-branch
cost by $1.9\times10^3\sigma^2$ at SNR 300 ($2.1\times10^4\sigma^2$ at SNR 1000), a likelihood ratio of about
$\E^{-935}$ or smaller. Complex ML selects the correct branch automatically.

\paragraph{5. The FID-sign test restores the bit.} The test of~\cite{cheng2019},
\begin{equation}
  \sgn\,\Re\sum_jS_{1j0}\overline{S_{2j0}}=\sgn(\sin\alpha_1\sin\alpha_2),
  \label{eq:fidsign}
\end{equation}
uses only the sign of the relative phase of the two FID images (it follows from
$a_0\propto v\propto\sin\alpha$ and the shared echo phases). In the magnitude fit it becomes a
constraint: $\Bone$ is confined to the interval between consecutive zeros of
$\sin(\Bone\alpha_i^{\rm nom})$ on which $\sin\alpha_1\sin\alpha_2$ has the measured sign; a start on the wrong
side is reflected across the nearest zero, where $|a_k|$ is symmetric, and every step is
projected back into the interval (\cref{alg:branch}).

\paragraph{6. What the test cannot do.} Branches with the same sign pattern (the complex
aliases, $\alpha_2>540^\circ$) are invisible to the test and to complex ML; the range bound excludes
them. Near $\alpha_2^{\rm eff}=360^\circ$ the scan-2 FID vanishes and the sign is unreliable, but the two
branches merge there, so a wrong decision costs little. Both the test and complex ML need the
relative phase of the two scans to be preserved: the same receiver phase and coil combination
and no motion or $B_0$ drift between the sequential acquisitions. With an unknown phase per
scan, complex ML would lose the relative sign and the alias would reappear for complex data
too.

\subsection{Comparison on identical data}

\Cref{tab:mag} compares the analytic inverse, magnitude least squares without and with the
FID-sign constraint, and complex ML on the noise realisations of \cref{tab:mc}.

\begin{table}[h]
\centering
\caption{Magnitude versus complex estimation on identical data (paper protocol, reference tissue, 1000 realisations). SD and mean of the log-estimates; bounds for complex data and for magnitude (Rician) data. Last column: voxels on the wrong side of $\alpha_2=360^\circ$ (``wrong'').}
\label{tab:mag}
\footnotesize
\setlength{\tabcolsep}{4pt}
\begin{tabular}{llccccc}
\toprule
SNR & & $\Tone$ SD & $\Tone$ mean & $\Bone$ SD & $\Ttwo$ SD & wrong\\
\midrule
300 & CRLB complex/Rician & 0.1350/0.1356 & & 0.0097/0.0098 & 0.0769/0.0771 & \\
 & Analytic inverse~\cite{cheng2019} & 0.9587 & +0.445 & 0.2695 & 0.5060 & --\\
 & Magnitude LS & 0.1909 & -0.030 & 0.0692 & 0.0743 & 47\\
 & Magnitude LS + FID sign & 0.1251 & -0.005 & 0.0094 & 0.0742 & 0\\
 & Complex ML & 0.1281 & -0.007 & 0.0097 & 0.0749 & --\\
\midrule
1000 & CRLB complex/Rician & 0.0405/0.0405 & & 0.0029/0.0029 & 0.0231/0.0231 & \\
 & Analytic inverse~\cite{cheng2019} & 0.2072 & +0.036 & 0.0657 & 0.1180 & --\\
 & Magnitude LS & 0.0457 & +0.000 & 0.0104 & 0.0231 & 3\\
 & Magnitude LS + FID sign & 0.0409 & +0.001 & 0.0029 & 0.0231 & 0\\
 & Complex ML & 0.0411 & +0.001 & 0.0029 & 0.0231 & --\\
\midrule
3000 & CRLB complex/Rician & 0.0135/0.0135 & & 0.0010/0.0010 & 0.0077/0.0077 & \\
 & Analytic inverse~\cite{cheng2019} & 0.0425 & +0.001 & 0.0046 & 0.0373 & --\\
 & Magnitude LS & 0.0136 & +0.000 & 0.0010 & 0.0076 & 0\\
 & Magnitude LS + FID sign & 0.0136 & +0.000 & 0.0010 & 0.0076 & 0\\
 & Complex ML & 0.0136 & +0.000 & 0.0010 & 0.0076 & --\\
\bottomrule
\end{tabular}
\end{table}


With the branch test, magnitude least squares has SD/CRLB within the Monte Carlo error of
complex ML at every SNR and no wrong-branch voxels; the Rician floor adds a small bias at SNR
300 ($-0.005$ in $\log\Tone$). Without it, wrong-branch voxels dominate the $\Bone$ error (SD up to
seven times the robust spread) and bias $\Tone$ by about $(\Bone)^{-2}$. The analytic inverse, which uses
the same magnitudes and the same FID sign, is 3--7 times worse for $\Tone$ and up to 28 times
for $\Bone$: since magnitudes carry essentially all the local information and a joint magnitude
fit extracts it, its loss is algorithmic.
